\appendix
\ssr{ПРИЛОЖЕНИЕ А}

К РПЗ работы прилагаются следующие материалы:
\begin{enumerate}
	\item Презентация к курсовой работе, состоящая из 12 слайдов.
	\item USB-флеш-накопитель с исходными кодами, электронными версиями РПЗ и презентации. Структура файлов на накопителе:
	\begin{verbatim}
\rpz
|
----\report.pdf // Итоговая РПЗ
----\prez // Презентация
----\source // Исходные коды
	\end{verbatim}
\end{enumerate}

\clearpage

\setcounter{section}{0}
\renewcommand{\thesection}{\Asbuk{section}}
\renewcommand{\thelstlisting}{\Asbuk{section}.\arabic{lstlisting}}
\setcounter{lstlisting}{0}
\section{Листинги кода}

\subsection{Функция расчета преломления луча}

\begin{lstlisting}[caption={Функция \texttt{refract\_dir} --- расчет направления преломленного луча}, label=lst:refract_dir, language=C++]
Vec3 refract_dir(const Vec3& dir, const Vec3& normal,
                 double n1, double n2, bool& total_internal) {
    Vec3 n = normal;
    double cos_i = -dot(n, dir);
    double eta = n1 / n2;
    double k = 1.0 - eta * eta * (1.0 - cos_i * cos_i);
    if (k < 0.0) {
        total_internal = true;
        return reflect_dir(dir, n);
    }
    total_internal = false;
    return dir * eta + n * (eta * cos_i - std::sqrt(k));
}
\end{lstlisting}

\subsection{Функция аппроксимации Френеля}

\begin{lstlisting}[caption={Функция \texttt{fresnel\_schlick} --- аппроксимация Френеля по формуле Шлика}, label=lst:fresnel_schlick, language=C++]
double fresnel_schlick(double cos_i, double n1, double n2) {
    double r0 = ((n1 - n2) / (n1 + n2));
    r0 = r0 * r0;
    
    double cos_i_clamped = std::max(0.0, std::min(1.0, std::abs(cos_i)));
    double one_minus_cos = 1.0 - cos_i_clamped;
    double power5 = one_minus_cos * one_minus_cos * one_minus_cos * one_minus_cos * one_minus_cos;
    return r0 + (1.0 - r0) * power5;
}
\end{lstlisting}

\subsection{Обработка преломления в колпачке}

\begin{lstlisting}[caption={Метод \texttt{Shield::process} --- обработка преломления луча в защитном колпачке}, label=lst:shield_process, language=C++]
std::optional<Ray> Shield::process(const Ray& incoming, std::mt19937& rng) const {
    Vec3 to_inner = inner_center_ - incoming.origin;
    double proj = dot(to_inner, incoming.direction);
    double dist_sq = dot(to_inner, to_inner) - proj * proj;
    double inner_sq = inner_radius_ * inner_radius_;
    if (dist_sq > inner_sq) {
        return incoming;
    }

    double thc = std::sqrt(inner_sq - dist_sq);
    double t0 = proj - thc;
    if (t0 < 1e-6) {
        t0 = proj + thc;
        if (t0 < 1e-6) {
            return incoming;
        }
    }

    Vec3 entry = incoming.at(t0);
    Vec3 entry_normal = normalize(entry - inner_center_);
    bool total = false;
    Ray inside = refract(incoming, entry, entry_normal, 1.0, config_.refractive_index, total);
    if (total) {
        return inside;
    }

    Vec3 to_outer = center_ - inside.origin;
    double proj_outer = dot(to_outer, inside.direction);
    double dist_outer_sq = dot(to_outer, to_outer) - proj_outer * proj_outer;
    double radius_sq = config_.radius * config_.radius;
    if (dist_outer_sq > radius_sq) {
        return inside;
    }
    double thc_outer = std::sqrt(radius_sq - dist_outer_sq);
    double t_exit = proj_outer + thc_outer;
    Vec3 exit_point = inside.at(t_exit);
    Vec3 exit_normal = normalize(exit_point - center_);
    Ray outside = refract(inside, exit_point, exit_normal,
                          config_.refractive_index, 1.0, total);
    if (total) {
        Vec3 dir = normalize(inside.direction);
        double cos_i = -dot(exit_normal, dir);
        double fresnel_r = fresnel_schlick(cos_i, config_.refractive_index, 1.0);
        double reflected_intensity = inside.intensity * fresnel_r;
        outside = Ray{exit_point + exit_normal * 1e-4,
                      reflect_dir(inside.direction, exit_normal),
                      reflected_intensity, inside.color};
    }

    double absorption = std::max(0.0, 1.0 - config_.transparency);
    outside.intensity *= (1.0 - absorption);
    outside.color = outside.color * config_.color;

    if (outside.intensity < 1e-5) {
        return std::nullopt;
    }

    if (config_.scatter_factor > 0.0) {
        std::normal_distribution<double> dist(0.0, config_.scatter_factor);
        Vec3 jitter(dist(rng), dist(rng), dist(rng));
        Vec3 new_dir = normalize(outside.direction + jitter);
        outside.direction = new_dir;
    }

    return outside;
}
\end{lstlisting}

\subsection{Ядро трассировки лучей}

\begin{lstlisting}[caption={Метод \texttt{Renderer::trace\_ray} --- ядро алгоритма трассировки лучей}, label=lst:trace_ray, language=C++]
void Renderer::trace_ray(const Ray& ray,
                         const LightConfig& source,
                         const std::vector<LightConfig>& lights,
                         const WallConfig& wall,
                         const std::vector<SphereConfig>& spheres,
                         const RenderSettings& settings,
                         ImageBuffer& wall_buffer,
                         std::vector<ImageBuffer>& sphere_buffers) {
    if (ray.intensity < 1e-5) {
        return;
    }

    double closest_t = 1e10;
    bool hit_object = false;
    int hit_type = 0;
    int hit_index = 0;
    Vec3 hit_point;
    Vec3 hit_normal;

    for (size_t i = 0; i < spheres.size(); ++i) {
        Vec3 point, normal;
        double t = 0.0;
        if (intersect_sphere(ray, spheres[i], point, normal, t)) {
            if (t < closest_t && t > 1e-6) {
                closest_t = t;
                hit_object = true;
                hit_type = 1;
                hit_index = static_cast<int>(i);
                hit_point = point;
                hit_normal = normal;
            }
        }
    }

    Vec3 wall_point;
    double wall_t = 0.0;
    bool hit_wall = intersect_wall(ray, wall, wall_point, wall_t);
    if (hit_wall && wall_t < closest_t && wall_t > 1e-6) {
        closest_t = wall_t;
        hit_object = true;
        hit_type = 0;
        hit_point = wall_point;
        hit_normal = WALL_NORMAL;
    }

    if (!hit_object) {
        return;
    }

    double lambert = std::max(0.0, dot(hit_normal, ray.direction * -1.0));
    double final_intensity = ray.intensity * lambert;

    if (hit_type == 0) {
        double norm_x = (hit_point.x + wall.width * 0.5) / wall.width;
        int pixel_x = static_cast<int>(norm_x * wall.resolution_x);
        double norm_y = (hit_point.y + wall.height * 0.5) / wall.height;
        int pixel_y = static_cast<int>((1.0 - norm_y) * wall.resolution_y);
        wall_buffer.add_pixel(pixel_x, pixel_y, ray.color, final_intensity);
    } else if (hit_type == 1) {
        const auto& sphere = spheres[hit_index];
        double u, v;
        sphere_point_to_uv(hit_point, sphere, u, v);
        int pixel_x = static_cast<int>(u * sphere.resolution);
        int pixel_y = static_cast<int>((1.0 - v) * sphere.resolution);
        if (hit_index < static_cast<int>(sphere_buffers.size())) {
            sphere_buffers[hit_index].add_pixel(pixel_x, pixel_y, ray.color, final_intensity);
        }
    }
}
\end{lstlisting}

\clearpage

\ssr{ПРИЛОЖЕНИЕ Б}

Презентация.

